\section{Sistematika Penulisan}
\label{sec:sistematika-penulisan}

Tugas akhir ini disusun dalam lima bab dengan sistematika sebagai berikut:

\noindent\textbf{Bab I} memuat latar belakang masalah, rumusan masalah, tujuan, batasan, manfaat, kebaruan penelitian, serta sistematika penulisan.

\noindent\textbf{Bab II} membahas landasan teori yang diperlukan, meliputi persamaan medan Einstein dan solusi Schwarzschild, konsep horizon peristiwa pada ruang-waktu Schwarzschild dan Schwarzschild Anti-de Sitter, teori perturbasi medan skalar tak bermassa dan bermassa, mode kuasinormal beserta kondisi batasnya, serta metode Frobenius untuk mengonstruksi ansatz solusi pada masing-masing ruang-waktu.

\noindent\textbf{Bab III} memaparkan metode penelitian yang digunakan, meliputi tahapan-tahapan penelitian dan diagram alir penyelesaian masalah.

\noindent\textbf{Bab IV} menyajikan hasil utama penelitian, yaitu penurunan relasi rekursif tiga suku beserta metode pecahan berlanjut, spektrum frekuensi mode kuasinormal pada ruang-waktu Schwarzschild dan Schwarzschild Anti-de Sitter untuk perturbasi medan skalar tak bermassa dan bermassa, serta analisis kestabilan kedua ruang-waktu terhadap perturbasi tersebut.

\noindent\textbf{Bab V} merangkum kesimpulan dari hasil-hasil penelitian dan menyajikan saran untuk penelitian lanjutan.